\fancyhead[C]{\textit{Solutions}}

\textbf{Note}: The empty line under \textbf{Main Idea \#1} should be \textit{linear dependence}.

\section{Solutions to Possible/Impossible Questions}
	\begin{table}[!h]
		\centering
		\begin{tabular}{c@{\hspace{2.5em}}cp{0.72\textwidth}}
			possible & impossible & \\
			\specialrule{\lightrulewidth}{0.25em}{1.25em}
			\bubble & \dbubble & \(\{\vec{v}_1, \vec{v}_2, \vec{v}_3\}\) forms a linear independent set of vectors from \(\mathbb{R}^n\) but
						      \(\{\vec{v}_1, \vec{v}_1 + \vec{v}_2, \vec{v}_3\}\) forms a linearly dependent set. \\[0.75em]
			\dbubble & \bubble & Two vectors \(\vec{v}_1, \vec{v}_2 \in \mathbb{R}^n\) form a linearly dependent set but one of the vectors
						      is not in the span of the other vector.\\[0.75em]
			\bubble & \dbubble & A matrix \(A \in \mathbb{R}^{2\times3}\) where the columns of \(A\) form a linearly independent set.\\[0.75em]
			\bubble & \dbubble & A matrix \(A \in \mathbb{R}^{m \times n}\) whose columns form a linearly independent set, and there exists an infinite 
						      number of solutions to the consistent system \(A\vec{x} = \vec{b}\) for some \(\vec{b} \in \mathbb{R}^m\).\\[0.75em]
			\bubble & \dbubble & A linearly independent set \(\{\vec{v}_1, \vec{v}_2, \vec{v}_3\}\) where \(\{\vec{v}_1, \vec{v}_2\}\) forms a linearly
						      dependent set.\\[0.75em]
			\bottomrule
		\end{tabular}
	\end{table}


\section{Solutions to Short Answers Questions}
\begin{enumerate}
	\item \(h = 3\). To see this, row reduce the following matrix to get 
		\[\begin{bmatrix}-2 & 1 & 7 \\ 6 & h & -3 \\ 1 & 2 & 4\end{bmatrix} \sim \cdots \sim \begin{bmatrix}1 & 0 & -2 \\ 0 & 1 & 3 \\ 0 & 0 & 9-3h\end{bmatrix}.\]
		Since the matrix has only 2 pivots, it spans a plane in \(\mathbb{R}^3\).
	\item Yes. Consider the case when \(\vec{b} = \vec{0}\). That means that \(A\vec{x} = \vec{0}\) only has the trivial solution as a solution. We immediately 
		then see that the columns of \(A\) must form a linearly independent set. (If it's still hard to see, write out the matrix-vector product as a linear 
		combination of the columns of \(A\),)
	\item An example could be \[2\begin{bmatrix}1 \\ 1 \\ 3\end{bmatrix} + 0\begin{bmatrix}0 \\ 1 \\ 0\end{bmatrix} + 1\begin{bmatrix}-2 \\ -2 \\ -6\end{bmatrix} = \vec{0}.\]
\end{enumerate}
